% Part: turing-machines
% Chapter: machines-computations
% Section: variants

\documentclass[../../../include/open-logic-section]{subfiles}

\begin{document}

\olfileid[fr]{tur}{mac}{var}
\olsection{Variantes des machines de Turing}

Il existe de nombreuses façons de définir les machines de Turing;
la nôtre n'en est qu'une. À certains égards, elle est plus
permissive que d'autres. Nous autorisons tout alphabet fini,
alors qu'une définition plus restrictive pourrait ne permettre
que deux symboles de ruban, $\TMstroke$ et~$\TMblank$.
Nous permettons à la machine d'écrire un symbole et de déplacer
sa tête au cours d'une même étape; d'autres définitions lui
permettent soit d'écrire, soit de se déplacer. Nous admettons
aussi qu'elle écrive sans déplacer la tête, tandis que
d'autres omettent le mouvement~$\TMstay$.

À d'autres égards, notre définition est plus restrictive.
Nous avons supposé que le ruban n'est infini que dans une
direction; d'autres définitions le rendent infini à gauche
comme à droite. On peut même autoriser plusieurs rubans
distincts, voire une grille infinie de cases. Nous représentons
les instructions de la machine par une fonction de transition;
d'autres définitions emploient une relation de transition,
qui peut offrir plusieurs instructions dans une même situation.

Cette dernière variante est particulièrement intéressante.
Dans notre définition, lorsque la machine est dans l'état~$q$
et lit le symbole~$\sigma$, la valeur
$\delta(q,\sigma)$ détermine le nouveau symbole, le nouvel
état et la position suivante de la tête. Mais si les
instructions forment une relation entre les paires
état-symbole courantes $\tuple{q,\sigma}$ et les triplets
nouvel état-symbole-direction
$\tuple{q',\sigma',D}$, l'action de la machine n'est
plus nécessairement unique : la relation peut contenir
à la fois $\tuple{q,\sigma,q',\sigma',D}$ et
$\tuple{q,\sigma,q'',\sigma'',D'}$.
Nous obtenons alors une machine de Turing
\emph{non déterministe}. De telles machines jouent
un rôle important en théorie de la complexité
algorithmique.

Les conventions d'arrêt varient également. Selon la nôtre,
la machine s'arrête lorsque sa fonction de transition n'est
pas définie; d'autres exigent qu'elle entre dans un état
d'arrêt désigné. Nous avons expliqué dans la
\olref[hal]{sec} que cette exigence ne change pas ce que
les machines de Turing peuvent calculer. Puisque nos
rubans ne sont infinis que dans une direction, une machine
peut recevoir une instruction qu'elle ne saurait exécuter
littéralement : déplacer sa tête vers la gauche alors qu'elle
lit la case la plus à gauche. Notre définition la laisse
sur place au lieu de la faire « tomber » hors du ruban,
mais on pourrait décider que la machine s'arrête alors.
Cette convention est également équivalente. Une première
idée, pour simuler une machine qui s'arrête lorsqu'elle
tente de quitter la case~$0$ vers la gauche, serait de
supprimer toute transition
$\delta(q,\TMendtape)=
\tuple{q',\sigma,\TMleft}$ : elle s'arrêterait alors
au lieu de tenter ce déplacement sur~$\TMendtape$.\footnote{
Précision éditoriale : nous explicitons aussi le cas où la case $0$ est écrasée, omis dans la réserve source. Cela ne suffit pas \emph{tout à fait} : rien n'interdit d'écrire
puis de lire~$\TMendtape$ dans une case autre que~$0$
(voir l'\olref[una]{ex:mover}), ni d'écraser le contenu
de la case~$0$. Une simulation correcte peut réserver
un marqueur distinct~$\TMendtape'$ à la case~$0$,
coder séparément le contenu de cette case et utiliser
un autre symbole pour les autres occurrences du marqueur.}

Les nombres, et donc les entrées et sorties des fonctions
calculées, peuvent aussi être représentés de différentes
manières. Nous employons l'écriture unaire, mais une
écriture binaire est également possible; elle nécessite
deux symboles en plus de~$\TMblank$ et de~$\TMendtape$.

Voici le fait remarquable : aucune de ces variantes
ne change l'ensemble des fonctions calculables par une
machine de Turing ; les ressources de temps ou d'espace peuvent, elles, changer. \emph{Une fonction calculable selon
l'une de ces définitions l'est aussi selon toutes les autres.}\footnote{Pour la variante non déterministe, on retient une valeur lorsqu'une branche finie la produit, et toutes les branches produisant une valeur doivent s'accorder sur celle-ci. L'exploration en largeur des branches finies retrouve cette valeur dès qu'elle existe ; l'existence de choix ne fournit pas un oracle.}

Nous ne démontrerons pas ce résultat en détail. Considérons
simplement les rubans infinis dans les deux directions ou
les machines à plusieurs rubans : ces variantes n'ajoutent
aucune puissance de calcul. En effet, une machine à un seul
ruban infini dans une direction peut les simuler. À l'aide
d'états et d'instructions supplémentaires, on peut, par
exemple, traduire le programme d'une machine à plusieurs
rubans ou à ruban bidirectionnel en celui d'une machine
à ruban unilatéral. Celle-ci peut affecter les cases paires
aux cases du premier ruban, ou aux positions « positives »
du ruban bidirectionnel, et les cases impaires aux cases
du second ruban, ou aux positions « négatives ».

\end{document}
